\documentclass[aspectratio=169,11pt]{beamer}
% Shared CMPSC 480 Beamer preamble.
\usepackage[T1]{fontenc}
\usepackage[utf8]{inputenc}
\usepackage{lmodern}
\usepackage{amsmath,amssymb,mathtools}
\usepackage{booktabs,array,tabularx}
\usepackage{xcolor}
\usepackage{listings}
\usepackage{tikz}
\usetikzlibrary{arrows.meta,positioning,shapes.geometric}

% Ignore only negligible TeX box diagnostics that are below visual tolerance.
% Larger layout defects still appear in the build log and in rendered QA.
\vfuzz=3pt
\hbadness=2000

\definecolor{courseink}{HTML}{111111}
\definecolor{coursegray}{HTML}{EDEDED}
\definecolor{courserule}{HTML}{B8BCC4}
\definecolor{courseblue}{HTML}{3D8DFF}
\definecolor{courselightblue}{HTML}{D0EDFA}
\definecolor{coursemuted}{HTML}{5C6470}
\definecolor{coursegreen}{HTML}{0B7A53}
\definecolor{coursered}{HTML}{B3261E}

\usetheme{default}
\usefonttheme{professionalfonts}
\setbeamertemplate{navigation symbols}{}
\setbeamersize{text margin left=0.65cm,text margin right=0.65cm}

\setbeamercolor{normal text}{fg=courseink,bg=white}
\setbeamercolor{frametitle}{fg=courseink,bg=white}
\setbeamercolor{title}{fg=courseink,bg=white}
\setbeamercolor{subtitle}{fg=coursemuted,bg=white}
\setbeamercolor{structure}{fg=courseblue}
\setbeamercolor{alerted text}{fg=coursered}
\setbeamercolor{block title}{fg=courseink,bg=courselightblue}
\setbeamercolor{block body}{fg=courseink,bg=coursegray}
\setbeamercolor{example text}{fg=coursegreen}

\setbeamerfont{title}{size=\fontsize{25}{29}\selectfont,series=\bfseries}
\setbeamerfont{subtitle}{size=\fontsize{13}{16}\selectfont}
\setbeamerfont{frametitle}{size=\fontsize{19}{22}\selectfont,series=\bfseries}
\setbeamerfont{normal text}{size=\fontsize{11}{14}\selectfont}
\setbeamerfont{footline}{size=\fontsize{7.5}{9}\selectfont}

\setbeamertemplate{frametitle}{%
  \nointerlineskip
  % Use content-driven height so a long assessment title can wrap without
  % being clipped by a fixed-height title bar.
  \begin{beamercolorbox}[wd=\paperwidth,sep=0.17cm,leftskip=0.48cm,rightskip=0.48cm]{frametitle}
    \insertframetitle
  \end{beamercolorbox}
  \vspace{-0.08cm}
  {\color{courserule}\rule{\textwidth}{0.35pt}}
}

\setbeamertemplate{footline}{%
  \leavevmode
  \begin{beamercolorbox}[wd=\paperwidth,ht=2.6ex,dp=1.1ex,leftskip=.65cm,rightskip=.65cm]{author in head/foot}
    \color{coursemuted}\insertshorttitle\hfill\insertframenumber/\inserttotalframenumber
  \end{beamercolorbox}
}

\setbeamertemplate{itemize item}{\color{courseblue}\small$\blacktriangleright$}
\setbeamertemplate{itemize subitem}{\color{courseblue}\scriptsize$\bullet$}
\setbeamertemplate{enumerate item}{\color{courseblue}\insertenumlabel.}

\lstdefinestyle{coursepython}{
  language=Python,
  basicstyle=\ttfamily\fontsize{8.5}{10}\selectfont,
  keywordstyle=\color{courseblue}\bfseries,
  commentstyle=\color{coursemuted}\itshape,
  stringstyle=\color{coursegreen},
  showstringspaces=false,
  columns=fullflexible,
  keepspaces=true,
  frame=single,
  rulecolor=\color{courserule},
  backgroundcolor=\color{coursegray},
  breaklines=true,
  tabsize=4,
  xleftmargin=0.15cm,
  xrightmargin=0.15cm,
  framexleftmargin=0.15cm,
  framexrightmargin=0.15cm
}
\lstset{style=coursepython}

\newcommand{\points}[1]{\hfill{\color{courseblue}\textbf{[#1 points]}}}
\newcommand{\deliverable}[1]{\textbf{\color{courseblue}#1}}
\newcommand{\code}[1]{\texttt{#1}}
\newcommand{\keyidea}[1]{%
  \begin{beamercolorbox}[rounded=false,sep=0.55em,wd=\linewidth]{block body}
    \textbf{Key idea.} #1
  \end{beamercolorbox}
}
\newcommand{\answerlabel}{\textbf{\color{coursegreen}Answer.}\ }
\newcommand{\gradinglabel}{\textbf{\color{courseblue}Grading guidance.}\ }

\newenvironment{questionblock}[1]
  {\begin{block}{#1}}
  {\end{block}}

\newenvironment{solutionblock}[1]
  {\begin{block}{#1}}
  {\end{block}}

\AtBeginSection[]{
  \begin{frame}
    \vfill
    {\usebeamerfont{title}\color{courseink}\insertsectionhead\par}
    \vspace{0.4cm}
    {\color{courseblue}\rule{0.32\paperwidth}{3pt}}
    \vfill
  \end{frame}
}


% CMPSC480_TOTAL_POINTS: 100
% CMPSC480_ITEM_IDS: 1,2,3,4,5,6,7
\title[Midterm Solutions]{CMPSC 480 Midterm Examination}
\subtitle{Cumulative assessment of Weeks 1-7}
\author{CMPSC 480 - Instructor Solution Deck}
\date{}

\begin{document}


\begin{frame}
  \titlepage
\vspace{-0.35cm}
\begin{center}
\color{coursered}\bfseries Instructor material - do not distribute with the question deck
\end{center}
\end{frame}

\begin{frame}{Solution overview}
\begin{columns}[T,onlytextwidth]
\begin{column}{0.62\textwidth}
\Large This instructor deck provides a complete matched solution and grading guidance.

\vspace{0.35cm}
\normalsize
Coverage: Modules A and B: agents, search, adversarial reasoning, MDPs, RL, approximation, and Bayesian inference
\end{column}
\begin{column}{0.33\textwidth}
\begin{block}{Points}100 total\end{block}
\begin{block}{Expected time}110 minutes\end{block}
\begin{block}{Submission}Individual examination PDF\end{block}
\end{column}
\end{columns}
\end{frame}

\begin{frame}{Instructor verification checklist}
\small
\begin{itemize}
  \item Work independently and show intermediate algebra, frontier state, or factor state.
  \item One double-sided handwritten letter-size reference sheet is permitted.
  \item A nonprogrammable calculator is permitted; networked devices and communication are prohibited.
  \item State every tie-breaking, terminal, reward, and zero-division convention you use.
  \item Answers receive credit for justified reasoning; an unsupported final number may receive little credit.
  \item Academic integrity rules of the course and institution apply throughout the examination.
\end{itemize}
\end{frame}

\begin{frame}{Solution 1.a - Agent and search formulation (3 points)}
\small
\textbf{Question focus.} Give one concrete element for each PEAS component.

\vspace{0.25cm}
\answerlabel Performance includes timely collision-free delivery and battery use; environment includes corridors, doors, people, and stations; actuators include motion and lift commands; sensors include localization, range sensing, battery, and task messages.
\end{frame}

\begin{frame}{Solution 1.b - Agent and search formulation (3 points)}
\small
\textbf{Question focus.} Define states, actions, transition model, goal test, and path cost for deterministic route planning.

\vspace{0.25cm}
\answerlabel A state contains position and other path-relevant discrete facts such as carried item or battery bucket; actions are legal moves or task operations; transition applies one action; goal checks delivery completion; cost can combine time and energy with nonnegative units.
\end{frame}

\begin{frame}{Solution 1.c - Agent and search formulation (4 points)}
\small
\textbf{Question focus.} Explain why rational behavior does not imply omniscience and name two defensive interface checks.

\vspace{0.25cm}
\answerlabel Rationality maximizes expected performance using the available percept history and knowledge, so uncertainty can still produce failure. Checks include immutable state keys, legal successor validation, positive costs for UCS/A-star, deterministic ties, and explicit failure.
\end{frame}

\begin{frame}{Solution 2.a - UCS and A-star (5 points)}
\small
\textbf{Question focus.} Trace UCS until a goal is removed, including cheaper rediscovery.

\vspace{0.25cm}
\answerlabel Remove S(0), insert A(4),B(1); remove B(1), improve A to 2 and insert G(9); remove A(2), improve G to 5; remove G(5). Return S-B-A-G with cost 5.
\end{frame}

\begin{frame}{Solution 2.b - UCS and A-star (5 points)}
\small
\textbf{Question focus.} Trace A-star priorities and return the selected path.

\vspace{0.25cm}
\answerlabel After S, A has f=7 and B f=5; remove B, improved A has g=2,f=5 and G has f=9; remove A, improve G to g=f=5; remove G and return cost 5.
\end{frame}

\begin{frame}{Solution 2.c - UCS and A-star (5 points)}
\small
\textbf{Question focus.} Is the heuristic admissible and consistent? Justify from the graph.

\vspace{0.25cm}
\answerlabel True remaining costs are h-star(G)=0, h-star(A)=3, h-star(B)=4 via A, and h-star(S)=5 via B-A. The listed h equals these values, so it is admissible. Each edge satisfies h(n)<=c+h(n-prime), including S-B:5<=1+4 and B-A:4<=1+3, so it is consistent.
\end{frame}

\begin{frame}{Solution 3.a - Minimax and alpha-beta (5 points)}
\small
\textbf{Question focus.} Compute every minimax value and the selected action.

\vspace{0.25cm}
\answerlabel MIN(A)=4, MIN(B)=3, and MIN(C)=2. MAX chooses A with root value 4.
\end{frame}

\begin{frame}{Solution 3.b - Minimax and alpha-beta (5 points)}
\small
\textbf{Question focus.} Trace alpha-beta left to right and identify pruned leaves.

\vspace{0.25cm}
\answerlabel After A, root alpha=4. At B, first leaf gives beta=3, so alpha>=beta and leaf 8 is pruned. At C, first leaf 5 does not prune because beta=5; second leaf gives 2. The root still chooses A.
\end{frame}

\begin{frame}{Solution 3.c - Minimax and alpha-beta (5 points)}
\small
\textbf{Question focus.} Design a depth-cutoff evaluation with three features and state two tests.

\vspace{0.25cm}
\answerlabel A valid MAX-perspective weighted sum might use score difference, mobility difference, and threat or distance advantage with normalized scales. Test terminal states bypass the heuristic and alpha-beta returns the same value as plain minimax on generated small trees.
\end{frame}

\begin{frame}{Solution 4.a - MDP planning (5 points)}
\small
\textbf{Question focus.} Compute Q(s,a) under reward-on-transition semantics.

\vspace{0.25cm}
\answerlabel Q(s,a)=0.8[5+0]+0.2[0+0.9(2)]=4+0.36=4.36.
\end{frame}

\begin{frame}{Solution 4.b - MDP planning (5 points)}
\small
\textbf{Question focus.} Compute Q(s,b) and the greedy action.

\vspace{0.25cm}
\answerlabel Q(s,b)=3 because the transition terminates and has no bootstrap. Since 4.36>3, choose a.
\end{frame}

\begin{frame}{Solution 4.c - MDP planning (5 points)}
\small
\textbf{Question focus.} Explain why value iteration converges for a finite discounted MDP.

\vspace{0.25cm}
\answerlabel The Bellman optimality operator changes values by at most gamma times the previous max-norm difference: expectation and maximum do not enlarge the bound beyond gamma. Since gamma<1, it is a contraction with a unique fixed point.
\end{frame}

\begin{frame}{Solution 5.a - Q-learning and approximation (5 points)}
\small
\textbf{Question focus.} Compute current Q, target, TD error, and new weights.

\vspace{0.25cm}
\answerlabel Current Q=0.5-1=-0.5. Target=2+0.5(3)=3.5; delta=4. Increment=0.2(4)[1,2]=[0.8,1.6], so new w=[1.3,1.1].
\end{frame}

\begin{frame}{Solution 5.b - Q-learning and approximation (5 points)}
\small
\textbf{Question focus.} Recompute the update if the transition is terminal.

\vspace{0.25cm}
\answerlabel Target=2, delta=2-(-0.5)=2.5, increment=0.2(2.5)[1,2]=[0.5,1], and new w=[1,0.5].
\end{frame}

\begin{frame}{Solution 5.c - Q-learning and approximation (5 points)}
\small
\textbf{Question focus.} Name the deadly triad and two observable instability signals.

\vspace{0.25cm}
\answerlabel The triad is function approximation, bootstrapping, and off-policy learning. Signals include exploding weight or Q norms, increasing absolute TD-error quantiles, nonfinite values, oscillatory loss or return, and large cross-seed variance.
\end{frame}

\begin{frame}{Solution 6.a - Bayesian networks and inference (5 points)}
\small
\textbf{Question focus.} Compute P(R,S,W) for all three true.

\vspace{0.25cm}
\answerlabel Use the factorization P(R)P(S|R)P(W|R,S)=0.2(0.1)(0.99)=0.0198.
\end{frame}

\begin{frame}{Solution 6.b - Bayesian networks and inference (5 points)}
\small
\textbf{Question focus.} Write an exact expression for P(R|W) without evaluating missing CPT rows.

\vspace{0.25cm}
\answerlabel P(R|W) is proportional to P(R) sum\_s P(s|R)P(W|R,s); compute that mass for R=true and false using their CPT rows, then normalize the two masses.
\end{frame}

\begin{frame}{Solution 6.c - Bayesian networks and inference (5 points)}
\small
\textbf{Question focus.} If factors f(R,S) and g(S,W) are multiplied and S is summed out, what is the result scope and why can order matter?

\vspace{0.25cm}
\answerlabel The product scope is \{R,S,W\}; summing S yields a factor over \{R,W\}. Different elimination orders can create intermediate factors over more variables and therefore require exponentially more space and arithmetic.
\end{frame}

\begin{frame}{Solution 7.a - Integrated engineering scenario (5 points)}
\small
\textbf{Question focus.} Specify the belief-to-policy interface and three invariants.

\vspace{0.25cm}
\answerlabel Pass a fixed-order normalized finite belief vector, optionally with observable state features, never the hidden truth. Invariants include correct dimension, nonnegative entries summing to one, and legal-action-only output.
\end{frame}

\begin{frame}{Solution 7.b - Integrated engineering scenario (5 points)}
\small
\textbf{Question focus.} Design a controlled ablation showing whether belief tracking helps.

\vspace{0.25cm}
\answerlabel Compare the full agent with a fixed-prior or observation-only agent using identical training steps, seeds, environment cases, and epsilon-zero evaluation; report return, success, hazard rate, and variability.
\end{frame}

\begin{frame}{Solution 7.c - Integrated engineering scenario (5 points)}
\small
\textbf{Question focus.} State a response to impossible evidence and a posttraining safety gate.

\vspace{0.25cm}
\answerlabel Raise or apply a declared fallback posterior when normalization is zero. Before acceptance, require belief oracle tests, terminal-update tests, finite Q/weight limits, and hazard-rate performance across noise levels.
\end{frame}

\begin{frame}{Edge-case reasoning and expected behavior}
\small
\begin{itemize}
  \item Return an empty path.
  \item Use stable action order.
  \item Ignore the bootstrap.
  \item Raise or use the declared fallback.
  \item Choose a lower-growth elimination order and report maximum factor scope.
\end{itemize}
\end{frame}

\begin{frame}{Grading rubric - part 1}
\scriptsize
\begin{tabularx}{\textwidth}{>{\bfseries}p{0.28\textwidth} p{0.08\textwidth} X}
\toprule
Category & Points & Full-credit evidence \\
\midrule
Agent and search formulation & 10 & Components are precise and connected to observable behavior. \\
UCS and A-star & 15 & Costs, updates, and guarantee conditions are correct. \\
Minimax and alpha-beta & 15 & Values, bounds, pruning, and perspective are correct. \\
MDP planning & 15 & Expectations, terminal convention, and contraction logic are sound. \\
\bottomrule
\end{tabularx}
\end{frame}

\begin{frame}{Grading rubric - part 2}
\scriptsize
\begin{tabularx}{\textwidth}{>{\bfseries}p{0.28\textwidth} p{0.08\textwidth} X}
\toprule
Category & Points & Full-credit evidence \\
\midrule
Q-learning and approximation & 15 & Arithmetic and deadly-triad reasoning are precise. \\
Bayesian networks and inference & 15 & CPT selection, normalization plan, and factor scope are correct. \\
Integrated engineering scenario & 15 & The answer connects inference, learning, evaluation, and defensive controls. \\
\bottomrule
\end{tabularx}
\end{frame}

\begin{frame}{Reflection exemplar}
\small
\answerlabel A full-credit response names the problem, states the assumption, and explains the exact trace, equation, or guarantee affected.

\vspace{0.25cm}
\gradinglabel Award credit for a specific claim supported by technical evidence from the submitted work.
\end{frame}

\end{document}
